% !TeX program = xelatex
\documentclass[11pt]{article}

\usepackage[a4paper,margin=1.7cm]{geometry}
\usepackage{fontspec}
\usepackage{polyglossia}
\setdefaultlanguage{hebrew}
\setotherlanguage{english}
\setmainfont{Times New Roman}
\newfontfamily\hebrewfont{Times New Roman}
\usepackage{amsmath,amssymb}
\usepackage{parskip}
\usepackage{tikz}
\usetikzlibrary{arrows.meta}

\pagenumbering{gobble}
\newcounter{prepq}

\newcommand{\q}[2]{\stepcounter{prepq}{\large\bfseries שאלה \arabic{prepq}. #1}\par\textbf{סמנו את האפשרות הנכונה היחידה.}\par #2\medskip}
\newcommand{\optbox}[1]{\fbox{\parbox[t][2.15cm][t]{0.43\textwidth}{\raggedright #1}}}
\newcommand{\mcfour}[4]{
  \begin{center}
  \begin{tabular}{@{}c@{\hspace{0.02\textwidth}}c@{}}
    \optbox{#1} & \optbox{#2} \\
    \optbox{#3} & \optbox{#4}
  \end{tabular}
  \end{center}
}

% Simple transition system diagram
\newcommand{\twostatets}{
\begin{tikzpicture}[scale=0.7]
  \node[circle, draw, line width=1.5pt, fill=gray!20, double] (s0) at (0,0) {$s_0$};
  \node[circle, draw, line width=1.5pt, fill=white] (s1) at (1.2,0) {$s_1$};
  \node[circle, draw, line width=1.5pt, fill=white] (s2) at (2.4,0) {$s_2$};
  \draw[-{Latex[length=2mm]}] (-0.7,0) -- (s0);
  \draw[-{Latex[length=2mm]}, thick] (s0) -- (s1);
  \draw[-{Latex[length=2mm]}, thick] (s1) -- (s2);
  \draw[-{Latex[length=2mm]}, thick, out=70, in=110, looseness=1.5] (s2) to (s0);
\end{tikzpicture}
}

% Simple cycle diagram
\newcommand{\cyclethree}{
\begin{tikzpicture}[scale=0.65]
  \node[circle, draw, line width=1.5pt, fill=gray!20, double] (s0) at (0,0) {$s_0$};
  \node[circle, draw, line width=1.5pt, fill=white] (s1) at (1.2,1) {$s_1$};
  \node[circle, draw, line width=1.5pt, fill=white] (s2) at (1.2,-1) {$s_2$};
  \draw[-{Latex[length=2mm]}] (-0.7,0) -- (s0);
  \draw[-{Latex[length=2mm]}, thick] (s0) -- (s1);
  \draw[-{Latex[length=2mm]}, thick] (s1) -- (s2);
  \draw[-{Latex[length=2mm]}, thick] (s2) -- (s0);
\end{tikzpicture}
}

% Product state space
\newcommand{\producttwo}{
\begin{tikzpicture}[scale=0.55]
  \node[circle, draw, line width=1.5pt, fill=gray!20, double] (pp) at (0,0) {$p_0,q_0$};
  \node[circle, draw, line width=1.5pt, fill=white] (p1q0) at (1.2,0) {$p_1,q_0$};
  \node[circle, draw, line width=1.5pt, fill=white] (p0q1) at (0,-1.2) {$p_0,q_1$};
  \node[circle, draw, line width=1.5pt, fill=white] (p1q1) at (1.2,-1.2) {$p_1,q_1$};
  \draw[-{Latex[length=2mm]}, thick] (pp) -- (p1q0);
  \draw[-{Latex[length=2mm]}, thick] (pp) -- (p0q1);
  \draw[-{Latex[length=2mm]}, thick] (p1q0) -- (p1q1);
  \draw[-{Latex[length=2mm]}, thick] (p0q1) -- (p1q1);
\end{tikzpicture}
}

% Loop on initial state
\newcommand{\loopinit}{
\begin{tikzpicture}[scale=0.7]
  \node[circle, draw, line width=1.5pt, fill=gray!20, double] (s0) at (0,0) {$s_0$};
  \node[circle, draw, line width=1.5pt, fill=white] (s1) at (1.2,0) {$s_1$};
  \draw[-{Latex[length=2mm]}] (-0.7,0) -- (s0);
  \draw[-{Latex[length=2mm]}, thick, out=100, in=80, looseness=2] (s0) to (s0);
  \draw[-{Latex[length=2mm]}, thick] (s0) -- (s1);
\end{tikzpicture}
}

\begin{document}

{\LARGE \textbf{50 שאלות הכנה לבוחן}}\\
המסמך כולל 50 שאלות אמריקאיות. בכל שאלה יש תשובה נכונה אחת בלבד.

\q{התאמת מערכת מעברים לגרף תוכנית}{
התאימו לגרף תוכנית שיש בו מצב התחלתי עם שני מעברים יוצאים לשני מצבים שונים.

\medskip\noindent\hfill\twostatets\hfill\mbox{}\medskip

\mcfour
  {\textbf{א.} למצב ההתחלתי אין אף מעבר יוצא.}
  {\textbf{ב.} למצב ההתחלתי יש בדיוק שני מעברים יוצאים לשני מצבים שונים.}
  {\textbf{ג.} אין מעבר מהמצב השני למצב הראשון.}
  {\textbf{ד.} למצב ההתחלתי יש מעבר אחד בלבד.}
}

\q{לולאה במצב התחלתי}{
סמנו את התיאור היחיד שמתאים למבנה שבו יש לולאה על המצב ההתחלתי וגם מעבר למצב נוסף.

\medskip\noindent\hfill\loopinit\hfill\mbox{}\medskip

\mcfour
  {\textbf{א.} מצב התחלתי ללא לולאה וללא מעבר החוצה.}
  {\textbf{ב.} שני מצבים עם לולאה רק במצב השני.}
  {\textbf{ג.} מצב התחלתי עם לולאה עצמית ומעבר למצב אחר.}
  {\textbf{ד.} שלושה מצבים בלי שום לולאות.}
}

\q{מחזור באורך 3}{
סמנו את מערכת המעברים היחידה שמתאימה למסלול $s_0\to s_1\to s_2\to s_0$.

\medskip\noindent\hfill\cyclethree\hfill\mbox{}\medskip

\mcfour
  {\textbf{א.} שני מצבים בלבד עם מעברים חוזרים.}
  {\textbf{ב.} שלושה מצבים עם מעברים יוצרים מחזור סגור.}
  {\textbf{ג.} מסלול באורך 3 ללא חזרה למצב ההתחלתי.}
  {\textbf{ד.} מצב יחיד עם שלוש לולאות נפרדות.}
}

\q{מצב סופי ללא יציאה}{
בגרף תוכנית מסוים המצב האחרון הוא מצב ללא מעברים יוצאים. סמנו את ההתאמה הנכונה.
\mcfour
  {\textbf{א.} המצב האחרון חייב להכיל לולאה עצמית לשלום עבודת התוכנית.}
  {\textbf{ב.} כל המצבים חייבים להיות התחלתיים כדי להבטיח סיום.}
  {\textbf{ג.} ייתכן שהמצב האחרון יהיה מצב סופי ללא שום מעבר יוצא.}
  {\textbf{ד.} צריך לסגור את הגרף בלולאה דרך המצב ההתחלתי.}
}

\q{התאמה למבנה קווי}{
סמנו את האפשרות היחידה שמתאימה למבנה קווי $s_0\to s_1\to s_2$ בלי חזרה לאחור.
\mcfour
  {\textbf{א.} מחזור באורך 2.}
  {\textbf{ב.} שלושה מצבים עם שני מעברים קדימה בלבד.}
  {\textbf{ג.} מצב יחיד עם שתי לולאות.}
  {\textbf{ד.} שני מצבים עם מעבר דו־כיווני.}
}

\q{שזירה ללא פעולות משותפות}{
יהיו שתי מערכות $A,B$ ללא פעולות משותפות. סמנו את הקביעה הנכונה.
\mcfour
  {\textbf{א.} כל צעד בהכרח עובר דרך שתי מערכות בו־זמנית.}
  {\textbf{ב.} מתקבלים צעדים בהם מערכת אחת מתחזקת בזמן שהשנייה מחכה.}
  {\textbf{ג.} כל צעד בתוצר מתאים לצעד בדיוק אחת מן המערכות.}
  {\textbf{ד.} המכפלה מצטמצמת למצב אחד בלבד מסיבות טכניות.}
}

\q{סנכרון על פעולה אחת}{
אם שתי המערכות יכולות לבצע את $a$ ו-$a$ שייכת לקבוצת הפעולות המשותפות, סמנו את הקביעה הנכונה.
\mcfour
  {\textbf{א.} ייתכן מעבר מסונכרן על $a$ אם שתי המערכות זמינות בו־זמנית לביצוע $a$.}
  {\textbf{ב.} $a$ תופיע כצעד של הראשונה בלבד ו־המערכת השנייה תחכה תמיד.}
  {\textbf{ג.} כאן אפשר רק צעד אינטרליבינג ממערכת אחת ללא סנכרון.}
  {\textbf{ד.} אם רק מערכת אחת יכולה לבצע $a$, גם אז זה נחשב צעד מסונכרן.}
}

\q{מספר מצבים במכפלה}{
למערכת $A$ יש 2 מצבים ולמערכת $B$ יש 3 מצבים. סמנו את המספר המקסימלי של מצבים במכפלה.
\mcfour
  {\textbf{א.} $5$}
  {\textbf{ב.} $6$}
  {\textbf{ג.} $2 + 3 = 5$ (מספר הצעדים האפשריים)}
  {\textbf{ד.} $4$ (אם הן מסונכרנות במלואן)}
}

\q{פעולה משותפת שאינה זמינה בתחילה}{
בכל אחת משתי המערכות הפעולה $a$ זמינה רק אחרי צעד אחד מקדים. סמנו את הקביעה הנכונה על המצב ההתחלתי.
\mcfour
  {\textbf{א.} במצב ההתחלתי אנחנו תמיד עורכים צעדי הכנה בלי סנכרון על $a$.}
  {\textbf{ב.} אפשר כל צעד מסונכרן אחרי שנבצע פעולות מקדימות בשתי המערכות.}
  {\textbf{ג.} המוצר חייב להיעצר במצב ההתחלתי כי $a$ אינה זמינה.}
  {\textbf{ד.} במצב ההתחלתי בהכרח מתרחש סנכרון על משהו אחר.}
}

\q{שתי פעולות משותפות}{
אם קבוצת הסנכרון היא $\{a,b\}$ ורק $a$ זמינה כעת בשתי המערכות, סמנו את הקביעה הנכונה.
\mcfour
  {\textbf{א.} שתי הפעולות חייבות להופיע יחד באותו צעד כדי להחשבות סנכרון.}
  {\textbf{ב.} קיים צעד מסונכרן על $a$ בלבד, כיוון שרק היא זמינה.}
  {\textbf{ג.} אנחנו חייבים לעצור כל עוד $b$ לא זמינה.}
  {\textbf{ד.} ניתן לבצע צעדי אינטרליבינג בלבד עד שגם $b$ תהיה זמינה.}
}

\q{אוגר שמתעדכן ב-$x$}{
נניח שאוגר יחיד מתעדכן לפי $r':=x$. סמנו את מספר המצבים המינימלי של מערכת המעברים המתארת את ערך האוגר בלבד.
\mcfour
  {\textbf{א.} $3$ מצבים (לערכים $0, 1$ ועוד).}
  {\textbf{ב.} צריך כמה מצבים כמו הערכים האפשריים של $x$.}
  {\textbf{ג.} $2$ מצבים לערכים בינאריים בלבד.}
  {\textbf{ד.} אי אפשר לתאר זאת אם $x$ יכול לקבל ערכים כמו משתנה אחר.}
}

\q{אוגר שמתאפס תמיד}{
אם בכל צעד מתקיים $r':=0$, סמנו את הקביעה הנכונה.
\mcfour
  {\textbf{א.} אחרי צעד אחד כל ריצה מגיעה למצב $0$ ונשארת בו.}
  {\textbf{ב.} המערכת מחליפה תמיד בין $0$ ולערך הקודם של $r$.}
  {\textbf{ג.} האוגר עובר דרך מצבים שונים לפני שמתייצב על $0$.}
  {\textbf{ד.} הערך הבא של האוגר אינו נקבע באופן דטרמיניסטי.}
}

\q{אוגר שמתהפך}{
אם בכל צעד מתקיים $r':=\neg r$, סמנו את התיאור הנכון.
\mcfour
  {\textbf{א.} מצב יחיד עם מעבר אחד לעצמו לוגית.}
  {\textbf{ב.} מסלול סופי ללא מעברים חוזרים.}
  {\textbf{ג.} מחזור באורך 2 בין הערכים $0$ ו-$1$.}
  {\textbf{ד.} מחזור באורך 3 בגלל משתני ביניים.}
}

\q{מעגל ללא זיכרון פנימי}{
אם הפלט תלוי רק בקלט הנוכחי ואין כלל זיכרון פנימי, סמנו את הקביעה הנכונה.
\mcfour
  {\textbf{א.} נדרשים לפחות 4 מצבים לייצוג אפילו לוגי פשוט.}
  {\textbf{ב.} המערכת יכולה להיות בעלת ריבוי מצבים בגלל ערכים שונים של הקלט.}
  {\textbf{ג.} מערכת בת מצב אחד מספיקה אם רק מודלים את ההתנהגות המיידית.}
  {\textbf{ד.} צריך מצב נוסף לשגיאות חובה.}
}

\q{שני ביטי זיכרון}{
למערכת יש שני ביטי זיכרון בלתי תלויים. סמנו את מספר המצבים הטבעי למודל המלא של הזיכרון.
\mcfour
  {\textbf{א.} $2$ (בחירה בין שני ערכים של bit)}
  {\textbf{ב.} $4$ (כל הצירופים האפשריים)}
  {\textbf{ג.} $8$ (אם נוסיף מצב שגיאה)}
  {\textbf{ד.} $16$ (תוך הוספת מצבי ביניים)}
}

\q{רצף של שלוש פקודות}{
נתון הקוד $x:=1;\ y:=0;\ z:=z+1$. סמנו את התיאור הנכון של גרף התוכנית.
\mcfour
  {\textbf{א.} תחילה קוד קווי של שלושה מעברים ברצף, ללא הסתעפויות.}
  {\textbf{ב.} הסתעפות בגלל שלוש השמות שונות.}
  {\textbf{ג.} שלוש לולאות עצמיות על אותו מצב.}
  {\textbf{ד.} הסתעפות אחרי הפקודה הראשונה בלבד.}
}

\q{פקודת if}{
לקוד $\mathbf{if}\ g_1\to x:=1\ \square\ g_2\to y:=1\ \mathbf{fi}$ סמנו את התיאור הנכון.
\mcfour
  {\textbf{א.} מצב כניסה עם שני מעברים יוצאים, כל אחד לענף.}
  {\textbf{ב.} לא ייתכנו ענפים - רק מעבר קווי אחד.}
  {\textbf{ג.} הגרף בהכרח מחזורי בגלל שני השומרים.}
  {\textbf{ד.} כל השומרים נמחקים ומוחלף כל אחד בלולאה.}
}

\q{פקודת do}{
לקוד $\mathbf{do}\ g\to x:=x+1\ \mathbf{od}$ סמנו את התיאור הנכון.
\mcfour
  {\textbf{א.} הגרף חייב להיות מסלול סופי.}
  {\textbf{ב.} אין שום אפשרות לחזרה לראש הלולאה.}
  {\textbf{ג.} הגרף כולל לולאה החוזרת לראש הלולאה.}
  {\textbf{ד.} כל השומרים נמחקים מן הגרף.}
}

\q{if בתוך רצף}{
לקוד $w:=0;\ \mathbf{if}\ g\to x:=1\ \square\ h\to y:=1\ \mathbf{fi};\ z:=2$ סמנו את התיאור הנכון.
\mcfour
  {\textbf{א.} תחילה קטע קווי, אחריו הסתעפות, ואז התכנסות חזרה לרצף.}
  {\textbf{ב.} כל הגרף הוא לולאה אחת בלבד.}
  {\textbf{ג.} חייבים להיות בדיוק שני מצבים.}
  {\textbf{ד.} אין מצב שבו שני ענפים מתכנסים.}
}

\q{skip בסוף רצף}{
אם הפקודה האחרונה בקוד היא $\mathrm{skip}$, סמנו את הקביעה הנכונה.
\mcfour
  {\textbf{א.} $\mathrm{skip}$ בהכרח מוחקת את שאר הפקודות.}
  {\textbf{ב.} $\mathrm{skip}$ מוסיפה בדרך כלל מעבר אחרון ללא שינוי מצב משתנים.}
  {\textbf{ג.} $\mathrm{skip}$ מכריחה לולאה חדשה.}
  {\textbf{ד.} $\mathrm{skip}$ מבטלת את מצב הסיום.}
}

\q{קריאה מערוץ ריק}{
בערוץ אסינכרוני ריק, מה נכון לגבי פעולה $a?x$?
\mcfour
  {\textbf{א.} הפעולה חסומה עד שתופיע הודעה בערוץ.}
  {\textbf{ב.} הפעולה תמיד מחזירה $0$.}
  {\textbf{ג.} הפעולה בוחרת ערך אקראי וממשיכה.}
  {\textbf{ד.} הפעולה יוצרת הודעה חדשה בערוץ.}
}

\q{קיבולת ערוץ 1}{
לערוץ $a$ קיבולת $1$. תהליך שולח פעמיים ברצף אל $a$ בלי שנעשית קריאה ביניים. סמנו את הקביעה הנכונה.
\mcfour
  {\textbf{א.} שתי השליחות תמיד מצליחות מייד.}
  {\textbf{ב.} השליחה השנייה עשויה להיחסם עד שתתבצע קריאה.}
  {\textbf{ג.} ההודעה הראשונה נמחקת אוטומטית.}
  {\textbf{ד.} הקיבולת משתנה אוטומטית ל-$2$.}
}

\q{מרוץ בין שתי שליחות}{
אם שני תהליכים שונים יכולים לשלוח הודעה לערוץ $b$ ותהליך שלישי מבצע קריאה יחידה מן הערוץ, סמנו את הקביעה הנכונה.
\mcfour
  {\textbf{א.} הערך שייקרא תלוי בסדר הריצה ויכול להגיע מכל אחד מן השולחים.}
  {\textbf{ב.} תמיד ייקרא רק הערך של השולח הראשון ברשימה.}
  {\textbf{ג.} הקריאה תחזיר את ממוצע שני הערכים.}
  {\textbf{ד.} הקריאה תחזיר תמיד את הערך הקטן יותר.}
}

\q{מצב גלובלי נגיש}{
שלושה תהליכים פועלים במקביל עם ערוצים סופיים. סמנו את התיאור הנכון של מצב גלובלי נגיש.
\mcfour
  {\textbf{א.} מצב גלובלי נגיש חייב לכלול תמיד ערוץ מלא.}
  {\textbf{ב.} מצב גלובלי נגיש חייב להשאיר לפחות תהליך אחד במצב ההתחלתי.}
  {\textbf{ג.} מצב גלובלי נגיש יכול לשלב מיקומים מקומיים שונים ותוכן ערוצים תואם.}
  {\textbf{ד.} מצב גלובלי נגיש אינו תלוי כלל בתוכן הערוצים.}
}

\q{סדר קריאות ושליחות}{
אם תהליך $R$ קורא תחילה מ-$a$ ורק אחר כך מ-$b$, סמנו את הקביעה הנכונה.
\mcfour
  {\textbf{א.} אפשר להגיע לקריאה מ-$b$ גם אם מעולם לא הושלמה קריאה מ-$a$.}
  {\textbf{ב.} הקריאה מ-$b$ מחייבת קודם מעבר מוצלח דרך הקריאה מ-$a$.}
  {\textbf{ג.} סדר הקריאות אינו משפיע על גרף המצבים המקומי של $R$.}
  {\textbf{ד.} הקריאה הראשונה חייבת תמיד להחזיר את ההודעה האחרונה שנשלחה למערכת כלשהי.}
}

\q{כלל גזירה חלופי לפיצול שומר והשמה}{
בכלל חלופי השומר וההשמה מפוצלים לשני מעברים. סמנו את הקביעה הנכונה.
\mcfour
  {\textbf{א.} בדרך כלל יופיע מצב ביניים בין בדיקת השומר ובין ביצוע ההשמה.}
  {\textbf{ב.} כל מצבי היציאה נמחקים.}
  {\textbf{ג.} כל הלולאות בגרף נעלמות.}
  {\textbf{ד.} המבנה הגרפי חייב להישאר זהה לחלוטין.}
}

\q{סיום רצף בכלל חלופי}{
אם בכלל חלופי סיום של $stmt_2$ מחזיר ישירות ל-$stmt_1$, סמנו את הקביעה הנכונה.
\mcfour
  {\textbf{א.} לעולם לא ייתכנו מעברים חוזרים.}
  {\textbf{ב.} עלולים להופיע מעברים החוזרים לתחילת הרצף.}
  {\textbf{ג.} כל רצף יהפוך למצב יחיד.}
  {\textbf{ד.} המשמעות הסמנטית של הרצף נעלמת לגמרי.}
}

\q{לולאת do בכלל חלופי}{
אם בכלל חלופי כל ענף של $\mathbf{do}$ חוזר ישירות לראש הלולאה, סמנו את התיאור הנכון.
\mcfour
  {\textbf{א.} מתקבל גרף קווי בלי חזרות.}
  {\textbf{ב.} מתקבל גרף שבו לכל ענף יש חזרה לראש הלולאה.}
  {\textbf{ג.} מתקבל גרף בלי שומרים כלל.}
  {\textbf{ד.} לא ייתכן כלל מצב ביניים.}
}

\q{skip בכלל חלופי}{
אם בגרסה חלופית $\mathrm{skip}$ מחזיר לראש הלולאה במקום לסיים, סמנו את הקביעה הנכונה.
\mcfour
  {\textbf{א.} עשויה להיווצר לולאה חדשה שלא קיימת בהגדרה הרגילה.}
  {\textbf{ב.} כל פקודות ההשמה נעלמות.}
  {\textbf{ג.} מספר המצבים תמיד קטן יותר.}
  {\textbf{ד.} אין שום הבדל מבני לעומת ההגדרה הרגילה.}
}

\q{השוואת שתי סמנטיקות}{
משווים בין סמנטיקה רגילה לסמנטיקה חלופית של אותה תוכנית. סמנו את הסימן המובהק ביותר לכך שהסמנטיקה השתנתה.
\mcfour
  {\textbf{א.} רק שמות המצבים השתנו אבל המבנה זהה.}
  {\textbf{ב.} כל התוויות הוחלפו ב-$\tau$ אוטומטית.}
  {\textbf{ג.} מספר מצבי הביניים או כיווני המעברים השתנו.}
  {\textbf{ד.} כל מצב קיבל מספר עמוד חדש.}
}

\q{בקר לרמזורים}{
סמנו את תכונת הבקר היחידה שתומכת בבטיחות של \textenglish{one green at a time}.
\mcfour
  {\textbf{א.} הבקר אפשר פתיחה בו־זמנית של כמה רמזורים כל עוד כמות מינימלית מחוברת.}
  {\textbf{ב.} הבקר מאפשר רמזור אחד בלבד ירוק בכל רגע, ויחכה לסיום לפני בחירה חדשה.}
  {\textbf{ג.} הבקר מתעלם מפעולות סיום ויוציא אור חדש בשרירותיות.}
  {\textbf{ד.} הבקר מסנכרן רק רמזורים שיש להם זוגות יחידים.}
}

\q{הפרת בטיחות במערכת רמזורים}{
סמנו את התרחיש היחיד שמרמז על הפרת בטיחות.
\mcfour
  {\textbf{א.} רמזור אחד נשאר אדום לאורך זמן בגלל רמזור שני שתופס הכל.}
  {\textbf{ב.} שני רמזורים שונים נמצאים בו־זמנית במצב ירוק.}
  {\textbf{ג.} הבקר חזר למצב התחלתי אחרי סדרה של פעולות.}
  {\textbf{ד.} קיימת לולאה עצמית על מצב מקומי של רכיב.}
}

\q{אלפבית הבקר}{
אם פעולת התחלה של רמזור מסוים אינה שייכת לאלפבית של הבקר, סמנו את הקביעה הנכונה.
\mcfour
  {\textbf{א.} הרמזור בהכרח יקבל עדיפות גבוהה בחירה ראשונה.}
  {\textbf{ב.} תכונת הבטיחות תישבר מייד ללא תנאי.}
  {\textbf{ג.} ייתכן שהרמזור כלל לא יוכל להתחיל לעולם.}
  {\textbf{ד.} כל שאר הרמזורים נעצרים אוטומטית.}
}

\q{אי־דטרמיניזם ובטיחות}{
אם הבקר בוחר אי־דטרמיניסטית בין $A_1$ ו-$A_2$, אבל אחרי כל התחלה הוא מחייב סיום מתאים לפני בחירה חדשה, סמנו את הקביעה הנכונה.
\mcfour
  {\textbf{א.} הבטיחות עדיין יכולה להישמר כל עוד סיום תמיד מתבצע.}
  {\textbf{ב.} הבטיחות בהכרח מופרעת בגלל בחירה אי־דטרמיניסטית.}
  {\textbf{ג.} המערכת הופכת לדטרמיניסטית אחרי מספר צעדים.}
  {\textbf{ד.} שני הרמזורים חייבים להיות ירוקים יחד לפחות פעם אחת.}
}

\q{השפעת קבוצת פעולות חלקית}{
אם קבוצת הפעולות המסונכרנות מכסה רק שני רמזורים ראשונים, סמנו את הקביעה הנכונה לגבי הרמזור השלישי.
\mcfour
  {\textbf{א.} הוא בהכרח יופעל ראשון.}
  {\textbf{ב.} הוא יכול להפוך לבקר הראשי.}
  {\textbf{ג.} ייתכן שהוא לא יקבל כלל אפשרות התחלה.}
  {\textbf{ד.} הוא חייב לבצע שתי התחלויות רצופות.}
}

\q{Neither safe nor live}{
סמנו את המצב היחיד שמתאים להגדרה \textenglish{neither safe nor live}.
\mcfour
  {\textbf{א.} המערכת בטוחה וחיה בכל ריצה.}
  {\textbf{ב.} יש הפרת בטיחות וגם קיימת ריצה של רעב או חסימה.}
  {\textbf{ג.} יש רק בעיית חיות אבל לעולם אין הפרת בטיחות.}
  {\textbf{ד.} יש רק הפרת בטיחות אבל כל הרכיבים מתקדמים תמיד.}
}

\q{חיות ללא בטיחות}{
סמנו את התיאור היחיד שמתאים למצב שבו חיות מתקיימת אבל בטיחות נכשלת.
\mcfour
  {\textbf{א.} כל רמזור מקבל ירוק אינסוף פעמים, אך לפעמים שני רמזורים ירוקים יחד.}
  {\textbf{ב.} תמיד יש חסימה אחרי מספר צעדים סופי.}
  {\textbf{ג.} לעולם אין מצב ירוק.}
  {\textbf{ד.} רמזור אחד בודד נשאר אדום לנצח וכל השאר כבויים.}
}

\q{בטיחות ללא חיות}{
סמנו את התיאור היחיד שמתאים למצב שבו בטיחות מתקיימת אבל חיות נכשלת.
\mcfour
  {\textbf{א.} בכל ריצה שני רמזורים ירוקים יחד.}
  {\textbf{ב.} לעולם לא שניים ירוקים יחד, אך ייתכן שרמזור אחד יורעב לנצח.}
  {\textbf{ג.} כל הרמזורים מחליפים צבע בהוגנות מלאה.}
  {\textbf{ד.} אין כלל מצבים במערכת.}
}

\q{Deadlock והשפעתו על חיות}{
סמנו את הקביעה הנכונה על \textenglish{deadlock} במערכת רמזורים.
\mcfour
  {\textbf{א.} הוא יכול לפגוע בחיות משום שמאותו מצב אין התקדמות נוספת.}
  {\textbf{ב.} הוא מחזק אוטומטית את תכונת הבטיחות.}
  {\textbf{ג.} הוא מבטיח שכל הרמזורים יקבלו ירוק.}
  {\textbf{ד.} הוא אפשרי רק אם כל הערוצים ריקים.}
}

\q{Scheduler לא הוגן}{
סמנו את הקביעה הנכונה על scheduler לא הוגן במערכת עם בחירה חוזרת בין שני רכיבים.
\mcfour
  {\textbf{א.} הוגנות מובטחת מעצם האי־דטרמיניזם.}
  {\textbf{ב.} scheduler לא הוגן מוחק מצבים מן המוצר.}
  {\textbf{ג.} scheduler לא הוגן עלול להרעיב רכיב שזמין שוב ושוב.}
  {\textbf{ד.} scheduler לא הוגן מחייב הפרת בטיחות מיידית.}
}

\q{קבוצת $H$ עם זוג מסונכרן אחד}{
נתונות שתי מערכות עם פעולות $a,c$ ו-$b,d$, ו-$H=\{\langle a,b\rangle\}$. סמנו את הקביעה הנכונה על הצעד הראשון.
\mcfour
  {\textbf{א.} מן המצב ההתחלתי אין שום צעד אפשרי.}
  {\textbf{ב.} מן המצב ההתחלתי ייתכנו צעדי אינטרליבינג עם $a$ בלבד או $b$ בלבד.}
  {\textbf{ג.} מן המצב ההתחלתי ייתכן רק צעד מסונכרן על $\langle a,b\rangle$.}
  {\textbf{ד.} ייתכנו צעדים נפרדים על $c$ ו-$d$ ביחד, ויתכן גם צעד מסונכרן על $\langle a,b\rangle$.}
}

\q{כאשר $H=\emptyset$}{
סמנו את הקביעה הנכונה כאשר $H=\emptyset$.
\mcfour
  {\textbf{א.} כל זוג פעולות חייב להיות מסונכרן אוטומטית.}
  {\textbf{ב.} נשארים רק צעדי אינטרליבינג, כאשר כל צעד של מערכת אחת בזה אחר זה.}
  {\textbf{ג.} המוצר מצטמצם למצב יחיד בגלל איבוד זוגות סנכרון.}
  {\textbf{ד.} כל התוויות במכפלה הופכות ל-$\tau$ אוטומטית.}
}

\q{כאשר רק $\langle c,d\rangle$ שייך ל-$H$}{
אם מן המצב ההתחלתי זמינות רק $a$ ו-$b$, אבל $H=\{\langle c,d\rangle\}$, סמנו את הקביעה הנכונה.
\mcfour
  {\textbf{א.} כבר מן ההתחלה חייב להופיע סנכרון על $\langle c,d\rangle$ כדי להתחיל.}
  {\textbf{ב.} לא ייתכן שום צעד במצב ההתחלתי כי $\langle c,d\rangle$ אינה זמינה.}
  {\textbf{ג.} במצב ההתחלתי ניתן לבצע צעדי אינטרליבינג על $a$ או $b$ בלבד.}
  {\textbf{ד.} כל צעדי האינטרליבינג נעלמים במצב זה ונוצר מצב סופי.}
}

\q{תווית של צעד מסונכרן}{
סמנו את הסימון הנכון היחיד לצעד מסונכרן בין פעולות $a,b$ בכללי שאלה זו.
\mcfour
  {\textbf{א.} $\langle a,b\rangle$}
  {\textbf{ב.} $a$ בלבד}
  {\textbf{ג.} $b$ בלבד}
  {\textbf{ד.} $\tau$ תמיד}
}

\q{המשך ריצה במצב $\langle p_1,q_1\rangle$}{
אם גם $\langle c,d\rangle$ שייך ל-$H$ ושתי הפעולות זמינות במצב $\langle p_1,q_1\rangle$, סמנו את הקביעה הנכונה.
\mcfour
  {\textbf{א.} לא ייתכנו עוד צעדים.}
  {\textbf{ב.} הצעד היחיד האפשרי הוא לולאה עצמית ללא תווית.}
  {\textbf{ג.} אפשר שיופיע במצב זה צעד מסונכרן נוסף על $\langle c,d\rangle$.}
  {\textbf{ד.} המערכת חייבת לחזור מיד למצב ההתחלתי.}
}

\q{מערכת שנייה בכיוון הפוך}{
בכלל חלופי שבו המערכת השנייה מתקדמת נגד כיוון החצים, נתון רק החץ $q_0\xrightarrow{b}q_1$. סמנו את הקביעה הנכונה על $b$ במצב $\langle p_0,q_0\rangle$.
\mcfour
  {\textbf{א.} בהכרח קיים צעד על $b^{rev}$ אל $\langle p_0,q_1\rangle$.}
  {\textbf{ב.} בהכרח קיים צעד רגיל על $b$.}
  {\textbf{ג.} נוצרה לולאה עצמית עם $b^{rev}$ בלי קשר לחצים.}
  {\textbf{ד.} מן הנתון הזה בלבד לא נובע צעד על $b$ מן המצב $\langle p_0,q_0\rangle$.}
}

\q{התקדמות לחילופין}{
אם מתחילים בתור של $TS_1$ ואז חובה להחליף תור אחרי כל צעד, סמנו את הקביעה הנכונה על שני הצעדים הראשונים.
\mcfour
  {\textbf{א.} שני הצעדים הראשונים יכולים להיות של $TS_1$.}
  {\textbf{ב.} הצעד הראשון של $TS_1$, והצעד השני של $TS_2$ אם הוא זמין.}
  {\textbf{ג.} שני הצעדים הראשונים חייבים להיות מסונכרנים.}
  {\textbf{ד.} אין כלל צעד ראשון מן המצב ההתחלתי.}
}

\q{סנכרון אחרי שני צעדים מכל צד}{
אם הסנכרון מוגדר רק אחרי זוג צעדים מכל מערכת, סמנו את התנאי ההכרחי.
\mcfour
  {\textbf{א.} בכל מערכת צריך להיות מסלול בן שני צעדים מן המצב הנוכחי.}
  {\textbf{ב.} מספיק שלכל מערכת תהיה פעולה אחת בלבד.}
  {\textbf{ג.} רק למערכת הראשונה נדרשים שני צעדים.}
  {\textbf{ד.} אין צורך כלל בקבוצת סנכרון.}
}

\q{כאשר אין כלל קבוצה $H$}{
אם במפורש לא מוגדרת קבוצה $H$, סמנו את הקביעה הנכונה.
\mcfour
  {\textbf{א.} נשארים רק צעדים של מערכת אחת בכל פעם, ללא סנכרון.}
  {\textbf{ב.} נשארים רק צעדים מסונכרנים כפולים.}
  {\textbf{ג.} כל מצב מקבל לולאה עצמית.}
  {\textbf{ד.} המערכת המשולבת ריקה.}
}

\q{צעד כפול לעומת אינטרליבינג}{
סמנו את ההבדל הנכון היחיד בין צעד מסונכרן כפול ובין אינטרליבינג רגיל.
\mcfour
  {\textbf{א.} באינטרליבינג רגיל תמיד קופצים מעל מצבי ביניים.}
  {\textbf{ב.} אין כל הבדל מבני בין שני סוגי הצעדים.}
  {\textbf{ג.} בצעד מסונכרן כפול מדלגים יחד מעל מצבי ביניים של שתי המערכות.}
  {\textbf{ד.} צעד מסונכרן כפול מוחק את כל המצבים הקודמים מן המודל.}
}

\end{document}
